\chapter{Lateral Internal Sphincterotomy}

\section*{Introduction \& Clinical Indications}
Lateral internal sphincterotomy (LIS) is a surgical procedure designed to alleviate the symptoms of chronic anal fissures by partially dividing the internal anal sphincter muscle. This condition is characterized by painful, linear tears in the anoderm, typically located in the posterior midline of the anal canal. The surgery aims to reduce the elevated resting anal pressure, which is a significant contributor to the development and persistence of these fissures. By decreasing the hypertonicity of the internal sphincter, LIS enhances blood flow to the affected area, promoting healing and providing substantial pain relief. This procedure is considered the gold standard for treating chronic anal fissures that have not responded to conservative management, including dietary modifications, topical treatments, and pharmacological therapies.

\begin{itemize}
  \item Internal Sphincterotomy
  \item Anal Fissure Surgery
  \item Lateral Sphincterotomy
  \item Open Lateral Internal Sphincterotomy (OLSI)
  \item Closed Lateral Internal Sphincterotomy (CLSI)
  \item Subcutaneous Lateral Internal Sphincterotomy
\end{itemize}

The surgical principle of lateral internal sphincterotomy revolves around the controlled division of the internal anal sphincter to alleviate hypertonicity, which is a primary factor in the pathogenesis of chronic anal fissures. The internal anal sphincter maintains a significant portion of the resting anal tone, and in patients with chronic fissures, this muscle often exhibits elevated resting pressure, leading to ischemia of the anoderm. This ischemia impedes healing and perpetuates a cycle of pain and muscle spasm.

The operative technique involves a precise incision through a portion of the internal anal sphincter, typically at the 3 o'clock or 9 o'clock position. This incision effectively reduces resting anal pressure by 30-50\%, improving local blood flow and promoting healing of the fissure. The technical goals of the procedure include:

\begin{itemize}
  \item Achieving a sustained reduction in resting anal pressure to facilitate fissure healing.
  \item Preserving the integrity of the external anal sphincter to prevent fecal incontinence.
  \item Minimizing trauma to surrounding tissues to reduce postoperative pain.
  \item Ensuring complete division of the targeted internal sphincter fibers without excessive resection.
  \item Avoiding injury to the dentate line and anal glands to prevent complications such as fistula formation.
\end{itemize}

The treatment of anal fissures has evolved significantly over time. Early approaches included caustic agents and cauterization, which often resulted in poor outcomes. The mid-20th century saw the emergence of techniques addressing anal sphincter spasm as a cause of fissures. In 1951, Eisenhammer introduced posterior internal sphincterotomy, which, while effective, had a higher incidence of complications.

A significant advancement occurred in 1967 when Sir Alan Parks pioneered the lateral internal sphincterotomy technique. This lateral approach minimized disruption to the anoderm and reduced the risk of complications associated with posterior techniques. Parks' method, performed either openly or subcutaneously, quickly gained acceptance due to its superior outcomes and lower morbidity. Subsequent refinements have focused on precise muscle division and minimal tissue dissection, solidifying LIS as the gold standard for chronic anal fissures.

Lateral internal sphincterotomy is indicated for patients with chronic anal fissures that have not responded to conservative management. The specific clinical indications include:

\begin{itemize}
  \item Persistent, severe anal pain, particularly during and after defecation, significantly impacting quality of life.
  \item Chronic anal fissures present for more than 6-8 weeks that have failed to heal with non-surgical treatments.
  \item Fissures unresponsive to topical pharmacological agents such as nitroglycerin or calcium channel blockers after a trial period.
  \item Recurrent anal fissures despite previous conservative treatments, indicating persistent internal sphincter hypertonicity.
  \item Presence of chronic changes such as a sentinel pile or hypertrophied anal papilla accompanying the fissure.
  \item Patients who have failed or are unsuitable for botulinum toxin injections due to contraindications or lack of efficacy.
  \item Objective evidence of high resting anal canal pressure on anorectal manometry.
  \item Fissures associated with significant bleeding unresponsive to conservative measures.
\end{itemize}

Lateral internal sphincterotomy is considered the primary surgical treatment for chronic anal fissures. It is not typically the first-line intervention for acute fissures, which often resolve with conservative measures. However, once an anal fissure becomes chronic, characterized by persistence beyond 6-8 weeks, LIS becomes the most effective treatment option. It is reserved for cases where conservative therapies have failed, providing superior healing rates and long-term efficacy compared to continued medical management.

\section*{Relevant Surgical Anatomy}
A comprehensive understanding of the anorectal anatomy is essential for the successful execution of lateral internal sphincterotomy. The anal canal, approximately 3-4 cm in length, extends from the anorectal ring to the anal verge. It is lined by columnar epithelium proximally, transitioning to specialized anoderm distally, which is highly innervated and sensitive. The dentate line serves as a critical anatomical landmark, marking the transition from rectal mucosa to anoderm and housing the anal crypts and papillae.

The anal continence mechanism relies on the internal and external anal sphincters. The internal anal sphincter (IAS) is a smooth muscle structure that contributes 70-85\% of the resting anal tone, measuring approximately 2.5-4 cm in length and 2-3 mm in thickness. The external anal sphincter (EAS) is a voluntary striated muscle that provides additional control over defecation. The intersphincteric groove, a palpable depression, separates the IAS and EAS and serves as a key surgical landmark. The arterial supply to the anal canal primarily comes from the inferior rectal arteries (branches of the internal pudendal artery), with venous drainage via the inferior rectal veins into the internal pudendal veins. The nerve supply to the EAS and anoderm is provided by the inferior rectal nerves (branches of the pudendal nerve), while the IAS is under autonomic control. Lymphatic drainage follows the inferior rectal vessels to the inguinal lymph nodes.

\section*{Preoperative Workup \& Preparation}
A thorough preoperative evaluation is essential for ensuring patient safety and optimizing surgical outcomes. This includes a detailed medical history focusing on bowel habits, previous anorectal surgeries, and current medications. A physical examination, including digital rectal examination and anoscopy, is performed to confirm the diagnosis of chronic anal fissure.

Routine laboratory tests, including complete blood count and coagulation profile, are ordered to assess overall health and bleeding risk. Specific preparations for the day of surgery include:

\begin{itemize}
  \item \textbf{NPO Status:} Patients are instructed to be NPO for 6-8 hours before surgery.
  \item \textbf{Bowel Preparation:} A mild laxative or cleansing enema may be prescribed to ensure an empty rectum.
  \item \textbf{Medication Review:} Current medications are reviewed, and anticoagulants may need to be adjusted.
  \item \textbf{Prophylactic Antibiotics:} Intravenous antibiotics may be administered prior to incision, especially in high-risk cases.
  \item \textbf{Informed Consent:} A detailed discussion of the procedure, risks, and expected outcomes is conducted, followed by obtaining informed consent.
\end{itemize}

While lateral internal sphincterotomy is generally safe, certain conditions may contraindicate its use or necessitate caution.

\textbf{Absolute Contraindications:}
\begin{itemize}
  \item Significant pre-existing fecal incontinence.
  \item Active inflammatory bowel disease affecting the anal canal.
  \item Active perianal sepsis, including abscess or symptomatic fistula.
  \item Uncorrected coagulopathy or severe bleeding disorders.
  \item Suspected or confirmed anal canal malignancy.
  \item Very low resting anal pressure on anorectal manometry.
\end{itemize}

\textbf{Relative Contraindications and Precautions:}
\begin{itemize}
  \item Advanced age, particularly in multiparous women.
  \item Prior anal surgery that may compromise sphincter function.
  \item Complex or atypical anal fissures requiring further investigation.
  \item Immunocompromised states that may impair wound healing.
  \item Pregnancy, where non-surgical management is preferred.
  \item History of obstetric anal sphincter injury.
\end{itemize}

\section*{Surgical Technique \& Procedure}
Lateral internal sphincterotomy is usually performed as an outpatient procedure under general or spinal anesthesia. Patient positioning is critical, typically in the prone jackknife or lithotomy position, with the perianal area prepped and draped in a sterile manner. A digital rectal examination is performed to assess anal tone and identify the intersphincteric groove.

The procedure can be executed using either an open or closed technique:

\begin{itemize}
  \item \textbf{Open Technique (Parks' Technique):}
    \begin{enumerate}
      \item A small transverse or oblique incision is made in the perianal skin at the 3 o'clock or 9 o'clock position.
      \item The anoderm and subcutaneous tissue are dissected to expose the intersphincteric groove.
      \item The internal anal sphincter is identified and a controlled incision is made through it, typically dividing 50-75\% of its thickness.
      \item Hemostasis is achieved using electrocautery.
      \item The skin incision may be left open to heal by secondary intention or loosely closed with absorbable sutures.
    \end{enumerate}

  \item \textbf{Closed (Subcutaneous) Technique (Notaras' Technique):}
    \begin{enumerate}
      \item A small stab incision is made in the perianal skin at the 3 o'clock or 9 o'clock position.
      \item A fine hemostat or scalpel is inserted into the intersphincteric groove to identify and elevate the internal anal sphincter.
      \item The internal sphincter is divided subcutaneously, ensuring the external sphincter is not injured.
      \item The small skin incision is typically left open to heal.
    \end{enumerate}
\end{itemize}

Both techniques aim to achieve effective sphincter division while minimizing trauma. The procedure usually takes 10-20 minutes, and no drains or implants are typically used.

\section*{Complications \& Risks}
Lateral internal sphincterotomy is associated with a low complication rate, but patients should be informed of potential risks.

\begin{itemize}
  \item \textbf{General Surgical Risks:}
    \begin{itemize}
      \item \textbf{Bleeding:} Minor bleeding is common; significant hemorrhage is rare.
      \item \textbf{Infection:} Wound infection may occur, requiring antibiotics or drainage.
      \item \textbf{Anesthesia Complications:} Risks include nausea, allergic reactions, and rare nerve damage.
      \item \textbf{Urinary Retention:} Temporary urinary retention may occur post-anesthesia.
    \end{itemize}
  \item \textbf{Procedure-Specific Risks:}
    \begin{itemize}
      \item \textbf{Fecal Incontinence:} The most significant complication, with varying degrees of severity.
      \item \textbf{Fissure Non-Healing or Recurrence:} A small percentage may fail to heal or recur.
      \item \textbf{Anal Fistula or Abscess Formation:} Rare but possible due to injury during dissection.
      \item \textbf{Persistent Pain:} Some patients may experience ongoing discomfort.
      \item \textbf{Hematoma Formation:} Blood collection under the skin may occur.
      \item \textbf{Keyhole Deformity:} Rarely occurs if excessive tissue is excised.
      \item \textbf{Perianal Edema or Swelling:} Common in the immediate postoperative period.
      \item \textbf{Mucosal Injury:} Accidental incision of the anal mucosa may occur.
    \end{itemize}
\end{itemize}

\section*{Postoperative Care \& Recovery}
Following lateral internal sphincterotomy, patients are monitored in the Post-Anesthesia Care Unit (PACU) until stable. Pain management is initiated with oral analgesics, and most patients are discharged the same day. 

In the first 24-48 hours, mild to moderate pain and serosanguinous drainage are common. Patients are encouraged to resume a high-fiber diet and use stool softeners to prevent straining. Sitz baths are recommended to promote hygiene and healing. Significant pain relief is typically experienced within the first week, with the fissure showing signs of healing within 2-3 weeks. Complete epithelialization usually occurs within 4-6 weeks, with most patients achieving symptom resolution and returning to normal bowel function.

During the procedure, the surgeon typically observes the chronic anal fissure as a deep, linear ulcer with indurated edges. The internal anal sphincter is palpated and confirmed to be hypertonic. Upon division, the muscle fibers retract, indicating successful release of tension. The absence of other perianal pathology is noted. If a sentinel pile or hypertrophied anal papilla is excised, these specimens are sent for pathological examination, confirming benign characteristics.

A successful postoperative course is characterized by rapid symptom resolution. Patients often experience immediate pain relief, with significant improvement within days. The first bowel movement post-surgery is typically less painful, especially with proper dietary management. The fissure begins to heal rapidly, with cessation of bleeding within the first week. The incision site heals over 4-6 weeks, with minimal scarring. Most patients return to normal bowel function without significant issues.

Postoperative care is vital for monitoring healing and preventing complications.

\begin{itemize}
  \item \textbf{Follow-up Visits:} Scheduled within 2-4 weeks to assess healing and address concerns.
  \item \textbf{Wound Care \& Hygiene:} Patients are instructed on proper hygiene and sitz baths to promote healing.
  \item \textbf{Activity Resumption:} Gradual return to normal activities is encouraged, with restrictions on strenuous activities for 2-4 weeks.
  \item \textbf{Dietary and Bowel Management:} Continued emphasis on a high-fiber diet and stool softeners is essential.
  \item \textbf{Warning Signs:} Patients are educated on signs of complications, including severe pain, fever, or significant bleeding.
\end{itemize}

Routine imaging or laboratory tests are generally not required unless complications arise.

\section*{Outcomes \& Prognosis}
Anal fissures are prevalent, with an estimated lifetime incidence of 1 in 10 individuals. The annual incidence is reported to be between 1 and 2 per 1,000 population. Chronic anal fissures, for which LIS is the definitive treatment, most commonly affect young to middle-aged adults, peaking between 30 and 50 years of age. There is no significant gender predilection, although a slight female predominance is noted, possibly linked to childbirth. The primary indication for LIS is chronic anal fissure refractory to conservative management, accounting for 30-50\% of cases. Mortality directly attributable to LIS is exceedingly rare, reflecting its safety profile.

Lateral internal sphincterotomy demonstrates excellent outcomes, with healing rates ranging from 90\% to 98\%. Most patients experience rapid pain relief and significant improvement in quality of life. The recurrence rate is low, typically less than 5\%. Prognostic factors for successful outcomes include proper patient selection and meticulous surgical technique. Long-term follow-up confirms sustained symptom resolution in most individuals, with a low incidence of persistent incontinence. The procedure's high efficacy and low complication profile solidify its position as the gold standard for chronic anal fissures.

\section*{References}
\begin{enumerate}
  \item MedlinePlus. Anal fissure surgery. U.S. National Library of Medicine; 2024. Available from: https://medlineplus.gov/ency/article/007384.htm.
  \item Schwartz's Principles of Surgery. 12th ed. McGraw-Hill Education; 2024. Chapter 30: Anorectal Disease.
  \item American Society of Colon and Rectal Surgeons (ASCRS). Clinical Practice Guidelines for the Management of Anal Fissures. Dis Colon Rectum. 2023;66(1):1-16.
  \item Corman ML. Corman's Colon and Rectal Surgery. 6th ed. Lippincott Williams \& Wilkins; 2013. Chapter 10: Anal Fissure.
\end{enumerate}

\clearpage
